\documentclass[10pt,a4paper]{article}
%% LuaLaTeX: на нём собран «Путеводитель Тьюринга», и только он умеет
%% microtype целиком — растяжение глифов и трекинг, которыми строка
%% подбирается сама. XeLaTeX эти два средства не поддерживает вовсе,
%% pdfLaTeX не годится: кириллических пакетов в среде сборки нет.
\usepackage{fontspec}
\usepackage{polyglossia}
\setmainlanguage{russian}
\setotherlanguage{english}
%% Начертания названы явно: без этого fontspec подставляет вместо курсива
%% ЖИРНЫЙ курсив, и все \textit в листке выходят жирными (проверено рендером).
%% Source Serif 4 — тот же шрифт, что выбран в html-конспектах проекта
%% (--serif в view.html). Файлы лежат рядом, в shrifty/, а не берутся
%% из системы: сборка должна давать один и тот же PDF в любой среде.
\setmainfont{SourceSerifPro}[
  Path           = shrifty/,
  Extension      = .ttf,
  UprightFont    = *_400Regular,
  ItalicFont     = *_400Regular_Italic,
  BoldFont       = *_700Bold,
  BoldItalicFont = *_700Bold_Italic]
\usepackage[left=11mm,right=13mm,top=18mm,bottom=14mm]{geometry}
\usepackage{amsmath,amssymb}
%% microtype с настройками из эталона владельца (_studio/konvejer/07-verstka/
%% etalon-tex/1.base-packages.sty): выступающая пунктуация и растяжение глифов
%% дают строке подобраться, и висящее на следующей строке слово втягивается
%% обратно — без этого приходится подгонять полосу руками.
%% Развилка по компилятору, а не выбор одного из двух. Целевой — LuaLaTeX:
%% на нём собран «Путеводитель Тьюринга» и только он умеет microtype целиком.
%% Но в песочнице, где собирается превью, luatex сломан (нет luatexbase.sty,
%% не грузятся шрифты Latin Modern, поставить неоткуда), и там работает
%% XeLaTeX — а он expansion и tracking не поддерживает и падает на них.
%% Так один и тот же .tex собирается в обеих средах, и в целевой — полностью.
\usepackage{iftex}
\ifLuaTeX
  \usepackage[protrusion=true,expansion=true,tracking=true,
              stretch=50,shrink=30,final]{microtype}
\else
  \usepackage[protrusion=true,final]{microtype}
\fi
\emergencystretch=1em
\linespread{0.98}
\frenchspacing
\setlength{\parindent}{0pt}
\pagestyle{empty}

%% Окно для отметки перед задачей — как в листках 179 школы: ученик отмечает
%% сданное сам. Ширина подобрана по образцу 06_Графы.pdf.
%% Окно для плюса. Размеры из архива 179 (\shag=1.7cm, высота 0.5\baselineskip),
%% ширина ужата: там окно на весь плюс с подписью, здесь хватает плюса.
\newlength{\shirinaokna}\setlength{\shirinaokna}{12mm}
\newlength{\vysotaokna}\setlength{\vysotaokna}{0.7\baselineskip}
\newcommand{\okno}{\framebox[\shirinaokna]{\rule{0pt}{\vysotaokna}}}
\newlength{\shirinanomera}\setlength{\shirinanomera}{9mm}
\newlength{\zazor}\setlength{\zazor}{0.6em}

%% Абзац задачи. Висячего отступа НЕТ намеренно: в листках 179 продолжение
%% строки начинается у левого поля, а не под текстом первой строки, и весь
%% текст листка стоит на одном левом крае. Окно и номер — обычный текст
%% в начале первой строки, а не колонка на поле. Висячий вариант пробовался
%% 03.09 и отвергнут владельцем: «почему „доползти до любой другой“ должно
%% начинаться там же, где „из спичек“».
\newcommand{\blok}[1]{\par\noindent #1\par}

%% Номер в боксе постоянной ширины: без него однозначные и двузначные номера
%% дают колонку разной ширины, и текст задач стоит на двух разных вертикалях.
\newcommand{\nomer}[2]{\makebox[\shirinanomera][l]{%
  \fbox{\rule{0pt}{\vysotaokna}\textbf{#1}$^{#2}$}}}
%% Термин при первом употреблении. Одна команда на весь листок: выделение
%% значит «у слова специальный смысл», а не «мне тут захотелось курсива».
\newcommand{\term}[1]{\textit{#1}}
%% Пустое место вместо окна — там, где плюс ставится не за задачу, а за пункты.
\newcommand{\pusto}{\hspace*{\shirinaokna}}
%% Знак письменной сдачи. Стандартный теховский кинжал, а не картинка пера:
%% pifont-перо в кегле листка не читается вовсе (проверено рендером 02.09).
%% Знак письменной сдачи — надстрочный, чтобы не съедал строку и не лез
%% в текст условия.
\newcommand{\pero}{\textsuperscript{\dag}\;}
\newcommand{\zvezdy}{%
  \par\medskip
  \centerline{\rule{0.3\linewidth}{0.4pt}\quad $*$\ \ $*$\ \ $*$\quad
              \rule{0.3\linewidth}{0.4pt}}%
  \par\medskip}

\begin{document}
\setbox0=\vbox{\hsize=\linewidth \begin{samepage}
\blok{\nomer{1}{\circ}\hspace{\zazor}Шесть компьютеров стоят в кабинете, и ни один ни с одним не соединён. Каждый день системный администратор соединяет кабелем какие-нибудь два. \term{Автономной зоной} назовём набор компьютеров, между любыми двумя из которых можно передать данные — напрямую или через другие компьютеры, — и такой, что ни с одним компьютером снаружи обмена нет.}
\blok{\okno\hspace{\zazor}\textbf{а)} Как может измениться число автономных зон за один день?}
\blok{\okno\hspace{\zazor}\textbf{б)} Какого наименьшего числа кабелей хватит, чтобы все шесть компьютеров оказались в одной зоне?}
\end{samepage}\par\vspace{3pt}}\typeout{MERA 1 0 \the\ht0}
\setbox0=\vbox{\hsize=\linewidth \begin{samepage}
\blok{\nomer{1}{\circ}\hspace{\zazor}Шесть компьютеров стоят в кабинете, и ни один ни с одним не соединён. Каждый день системный администратор соединяет кабелем какие-нибудь два. \term{Автономной зоной} назовём набор компьютеров, между любыми двумя из которых можно передать данные — напрямую или через другие компьютеры, — и такой, что ни с одним компьютером снаружи обмена нет.\hspace{\zazor}\okno\hspace{\zazor}\textbf{а)} Как может измениться число автономных зон за один день?\hspace{\zazor}\okno\hspace{\zazor}\textbf{б)} Какого наименьшего числа кабелей хватит, чтобы все шесть компьютеров оказались в одной зоне?}
\end{samepage}\par\vspace{3pt}}\typeout{MERA 1 1 \the\ht0}
\setbox0=\vbox{\hsize=\linewidth \begin{samepage}
\blok{\nomer{6}{\circ}\hspace{\zazor}В классе 20 человек. \term{Компанией} назовём набор ребят, любые двое из которых связаны цепочкой знакомств внутри этого набора и который ни с кем снаружи не знаком.}
\blok{\okno\hspace{\zazor}\textbf{а)} В классе ровно 16 пар знакомых. На какое наименьшее число компаний может разбиться класс?}
\blok{\okno\hspace{\zazor}\textbf{б)} Класс разбился ровно на четыре компании. Какое наименьшее число пар знакомых в нём?}
\end{samepage}\par\vspace{3pt}}\typeout{MERA 6 0 \the\ht0}
\setbox0=\vbox{\hsize=\linewidth \begin{samepage}
\blok{\nomer{6}{\circ}\hspace{\zazor}В классе 20 человек. \term{Компанией} назовём набор ребят, любые двое из которых связаны цепочкой знакомств внутри этого набора и который ни с кем снаружи не знаком.\hspace{\zazor}\okno\hspace{\zazor}\textbf{а)} В классе ровно 16 пар знакомых. На какое наименьшее число компаний может разбиться класс?\hspace{\zazor}\okno\hspace{\zazor}\textbf{б)} Класс разбился ровно на четыре компании. Какое наименьшее число пар знакомых в нём?}
\end{samepage}\par\vspace{3pt}}\typeout{MERA 6 1 \the\ht0}
\setbox0=\vbox{\hsize=\linewidth \zvezdy
\begin{samepage}
\blok{\nomer{7}{\circ}\hspace{\zazor}Дан граф с $n$ вершинами. Докажите, что}
\blok{\okno\hspace{\zazor}\textbf{а)} если граф связен, то в нём не менее $n-1$ ребра;}
\blok{\okno\hspace{\zazor}\textbf{б)} если граф распадается на $k$ компонент связности, то в нём не менее $n-k$ рёбер;}
\blok{\okno\hspace{\zazor}\textbf{в)} если в связном графе есть цикл, то любое ребро этого цикла можно удалить (концы остаются на месте), и граф останется связным.}
\end{samepage}\par\vspace{3pt}}\typeout{MERA 7 0 \the\ht0}
\setbox0=\vbox{\hsize=\linewidth \zvezdy
\begin{samepage}
\blok{\nomer{7}{\circ}\hspace{\zazor}Дан граф с $n$ вершинами. Докажите, что\hspace{\zazor}\okno\hspace{\zazor}\textbf{а)} если граф связен, то в нём не менее $n-1$ ребра;\hspace{\zazor}\okno\hspace{\zazor}\textbf{б)} если граф распадается на $k$ компонент связности, то в нём не менее $n-k$ рёбер;\hspace{\zazor}\okno\hspace{\zazor}\textbf{в)} если в связном графе есть цикл, то любое ребро этого цикла можно удалить (концы остаются на месте), и граф останется связным.}
\end{samepage}\par\vspace{3pt}}\typeout{MERA 7 1 \the\ht0}
\setbox0=\vbox{\hsize=\linewidth \begin{samepage}
\blok{\nomer{10}{\circ}\hspace{\zazor}\okno\hspace{\zazor}\textbf{а)} В связном графе 9 вершин и 14 рёбер. Какое наименьшее число рёбер надо удалить, чтобы в графе не осталось ни одного цикла?}
\blok{\okno\hspace{\zazor}\textbf{б)} Докажите, что из любого связного графа можно выкинуть часть рёбер так, чтобы получилось дерево с теми же вершинами. Такое дерево называют \term{остовом} исходного графа.}
\end{samepage}\par\vspace{3pt}}\typeout{MERA 10 0 \the\ht0}
\setbox0=\vbox{\hsize=\linewidth \begin{samepage}
\blok{\nomer{10}{\circ}\hspace{\zazor}\okno\hspace{\zazor}\textbf{а)} В связном графе 9 вершин и 14 рёбер. Какое наименьшее число рёбер надо удалить, чтобы в графе не осталось ни одного цикла?\hspace{\zazor}\okno\hspace{\zazor}\textbf{б)} Докажите, что из любого связного графа можно выкинуть часть рёбер так, чтобы получилось дерево с теми же вершинами. Такое дерево называют \term{остовом} исходного графа.}
\end{samepage}\par\vspace{3pt}}\typeout{MERA 10 1 \the\ht0}
\setbox0=\vbox{\hsize=\linewidth \begin{samepage}
\blok{\nomer{11}{\circ}\hspace{\zazor}Вершину степени 1 называют \term{висячей}, или \term{листом}.}
\blok{\okno\hspace{\zazor}\textbf{а)} Докажите, что в дереве больше чем с одной вершиной есть висячая вершина. \textit{(идите по рёбрам, каждый раз выбирая то, по которому ещё не шли)}}
\blok{\okno\hspace{\zazor}\textbf{б)} Докажите, что таких вершин хотя бы две.}
\blok{\okno\hspace{\zazor}\textbf{в)} Сколько висячих вершин может быть у дерева с $n$ вершинами?}
\end{samepage}\par\vspace{3pt}}\typeout{MERA 11 0 \the\ht0}
\setbox0=\vbox{\hsize=\linewidth \begin{samepage}
\blok{\nomer{11}{\circ}\hspace{\zazor}Вершину степени 1 называют \term{висячей}, или \term{листом}.\hspace{\zazor}\okno\hspace{\zazor}\textbf{а)} Докажите, что в дереве больше чем с одной вершиной есть висячая вершина. \textit{(идите по рёбрам, каждый раз выбирая то, по которому ещё не шли)}\hspace{\zazor}\okno\hspace{\zazor}\textbf{б)} Докажите, что таких вершин хотя бы две.\hspace{\zazor}\okno\hspace{\zazor}\textbf{в)} Сколько висячих вершин может быть у дерева с $n$ вершинами?}
\end{samepage}\par\vspace{3pt}}\typeout{MERA 11 1 \the\ht0}
\setbox0=\vbox{\hsize=\linewidth \begin{samepage}
\blok{\nomer{13}{\star}\hspace{\zazor}$N$-угольник разбит на треугольники несколькими диагоналями, не пересекающимися нигде, кроме вершин. Построим граф, соответствующий этому разбиению: отметим внутри каждого треугольника точку (это будут вершины графа), соединяя две точки ребром ровно в том случае, когда соответствующие точкам треугольники имеют общую сторону. Докажите, что}
\blok{\okno\hspace{\zazor}\textbf{а)} этот граф будет деревом;}
\blok{\okno\hspace{\zazor}\textbf{б)} хотя бы у двух треугольников разбиения две стороны совпадают со сторонами $N$-угольника (если треугольников больше одного).}
\end{samepage}\par\vspace{3pt}}\typeout{MERA 13 0 \the\ht0}
\setbox0=\vbox{\hsize=\linewidth \begin{samepage}
\blok{\nomer{13}{\star}\hspace{\zazor}$N$-угольник разбит на треугольники несколькими диагоналями, не пересекающимися нигде, кроме вершин. Построим граф, соответствующий этому разбиению: отметим внутри каждого треугольника точку (это будут вершины графа), соединяя две точки ребром ровно в том случае, когда соответствующие точкам треугольники имеют общую сторону. Докажите, что\hspace{\zazor}\okno\hspace{\zazor}\textbf{а)} этот граф будет деревом;\hspace{\zazor}\okno\hspace{\zazor}\textbf{б)} хотя бы у двух треугольников разбиения две стороны совпадают со сторонами $N$-угольника (если треугольников больше одного).}
\end{samepage}\par\vspace{3pt}}\typeout{MERA 13 1 \the\ht0}
\end{document}